\documentclass[11pt]{article}

\usepackage[margin=1.05in]{geometry}
\usepackage{amsmath,amssymb}
\usepackage{booktabs}
\usepackage{array}
\usepackage{verbatim}
\usepackage[colorlinks=true,linkcolor=blue,urlcolor=blue]{hyperref}

\newcommand{\bu}{\mathbf{u}}
\newcommand{\bq}{\mathbf{q}}
\newcommand{\buh}{\widehat{\mathbf{u}}}
\newcommand{\bn}{\mathbf{n}}
\newcommand{\bx}{\mathbf{x}}
\newcommand{\Minf}{M_\infty}
\newcommand{\Rey}{\mathrm{Re}}
\newcommand{\Pran}{\mathrm{Pr}}
\newcommand{\code}[1]{\texttt{#1}}

\title{An MFEM HDG Solver for Steady Hypersonic Viscous Flow:\\
Physics, Discretization, and Implementation\\
{\large \code{myapps/hdg\_navierstokes} --- replication of the Exasim
\code{nsmach8-baseline} run}}
\author{}
\date{July 2026}

\begin{document}
\maketitle

\begin{abstract}
This document describes the hybridizable discontinuous Galerkin (HDG) solver in
\code{myapps/hdg\_navierstokes}, an MFEM-based reimplementation of the Exasim
\code{nsmach8-baseline} case: steady laminar flow of a perfect gas at
$\Minf = 8.03$, $\Rey = 1.835\times 10^5$ over an isothermal half-cylinder.
The solver reproduces Exasim's discretization to residual-level parity
(the Exasim converged state evaluates to a residual of $3.3\times10^{-7}$ in
this code's assembler) and its converged fields to $\sim\!10^{-10}$ relative
$L^2$. We document the governing equations and their nondimensionalization
(including two deliberate idiosyncrasies of the reference run), the artificial
viscosity and the smoothed density/pressure regularization that the
continuation requires, the HDG formulation with static condensation, the
damped-Newton globalization and the four-stage artificial-viscosity
continuation, and the code layout, data formats, and verification gates.
\end{abstract}

\tableofcontents

%======================================================================
\section{Overview and scope}
%======================================================================

The target of this code is \emph{replication}, not merely simulation: it
reproduces a specific reference computation, the Exasim run archived at
\code{Exasim-apps/runs/nsmach8-baseline}, with the same mesh (identical
discrete geometry), the same physics model (transcribed term by term from the
generated model code), the same HDG discretization (spaces, stabilization,
quadrature), and the same nonlinear solver semantics (damped Newton with
Exasim's exact line-search and abort rules). The only deliberate deviation is
the linear solver (a sparse direct factorization instead of Exasim's
preconditioned GMRES), which changes iteration counts but not the Newton fixed
point.

Development was organized in four milestones. M1: physics/mesh/IO foundations
verified against the Exasim data structures. M2: the HDG core verified on
benign problems (manufactured solutions, freestream preservation, a Mach~4.5
sanity case). M3: the Mach~8 baseline replication from Exasim's restart data.
M4: a self-contained four-stage artificial-viscosity continuation from a
damped-freestream initial condition, requiring no Exasim state data. All four
milestone gates pass (\S\ref{sec:verification}).

%======================================================================
\section{Physical problem}
%======================================================================

The flow is two-dimensional, steady, laminar, compressible flow of a
calorically perfect gas over the front half of a circular cylinder:
\begin{center}
\begin{tabular}{ll}
\toprule
Mach number & $\Minf = 8.03$ \\
Reynolds number (based on cylinder radius) & $\Rey = 1.835\times 10^{5}$ \\
Prandtl number & $\Pran = 0.71$ \\
Specific-heat ratio & $\gamma = 1.4$ \\
Freestream static temperature & $T_{\mathrm{ref}} = 124.49~\mathrm{K}$ \\
Wall temperature (isothermal) & $T_{\mathrm{wall}} = 294.44~\mathrm{K}$ \\
Sutherland constant & $T_s = 110.4~\mathrm{K}$ \\
\bottomrule
\end{tabular}
\end{center}
The domain is a half-annulus: the inner boundary is the cylinder surface
(radius $1$ after scaling), the outer boundary is a smooth arc standing off
the bow shock, and the two straight segments on the $x=0$ plane are supersonic
outflow planes. A strong detached bow shock forms upstream of the cylinder
(standoff $\approx 0.4$ radii on the stagnation line); downstream of it the
flow compresses to $\rho \approx 5.6$ (inviscid Rankine--Hugoniot level), and
in the cold-wall thermal boundary layer density rises further (the converged
field reaches $\rho \approx 28$ at this Reynolds number). The steady solution
is found directly by Newton's method; there is no time stepping.

%======================================================================
\section{Governing equations}
%======================================================================

\subsection{Conservative form}

With conservative variables
$\bu = (\rho,\ \rho u,\ \rho v,\ \rho E)^{T}$, the steady compressible
Navier--Stokes equations are
\begin{equation}
\nabla\cdot\bigl(F_{\mathrm{inv}}(\bu) - F_{\mathrm{visc}}(\bu,\nabla\bu)
- F_{\mathrm{av}}(\nabla\bu)\bigr) = 0 ,
\label{eq:ns}
\end{equation}
\begin{equation}
F_{\mathrm{inv}} =
\begin{pmatrix}
\rho u & \rho v\\
\rho u^2 + p & \rho u v\\
\rho u v & \rho v^2 + p\\
\rho u H & \rho v H
\end{pmatrix},
\qquad
F_{\mathrm{visc}} =
\begin{pmatrix}
0 & 0\\
\tau_{xx} & \tau_{xy}\\
\tau_{xy} & \tau_{yy}\\
u\tau_{xx}+v\tau_{xy}+f_c\,\partial_x T &
u\tau_{xy}+v\tau_{yy}+f_c\,\partial_y T
\end{pmatrix},
\end{equation}
with $p = (\gamma-1)\bigl(\rho E - \tfrac12\rho(u^2+v^2)\bigr)$, total
enthalpy $H = E + p/\rho$, and Stokes-hypothesis stresses
\begin{equation}
\tau_{xx} = \mu\,\tfrac23\,(2u_x - v_y),\qquad
\tau_{xy} = \mu\,(u_y + v_x),\qquad
\tau_{yy} = \mu\,\tfrac23\,(2v_y - u_x).
\end{equation}
The ``temperature'' variable used throughout the model is the specific
internal energy,
\begin{equation}
T \;=\; \frac{p}{(\gamma-1)\,\rho} \;=\; e,
\end{equation}
and the heat conduction coefficient is $f_c = \mu\gamma/\Pran$, so that
$f_c \nabla T = (\mu\gamma/\Pran)\nabla e$ equals the Fourier flux
$ (\mu c_p / \Pran) \nabla T_{\mathrm{thermo}}$ in dimensional terms. This
convention (temperature $\equiv$ internal energy) matters when comparing wall
heat flux with other codes.

\subsection{Sutherland viscosity}

The dynamic viscosity follows Sutherland's law, evaluated algebraically from
the local state. Defining the nondimensional temperature ratio
\begin{equation}
T_r \;=\; \gamma \Minf^2\, \frac{p}{\rho}
\;=\; \frac{T}{T_\infty},
\qquad T_{\mathrm{phys}} = T_{\mathrm{ref}}\, T_r,
\end{equation}
the code computes
\begin{equation}
\mu \;=\; \frac{1}{\Rey}\; T_r^{3/2}\;
\frac{T_{\mathrm{ref}} + T_s}{T_{\mathrm{phys}} + T_s},
\qquad T_s = 110.4~\mathrm{K}.
\label{eq:sutherland}
\end{equation}
Note $T_r^{3/2}$: for $p<0$ this is the square root of a negative number ---
the origin of the robustness problem addressed by the regularization of
\S\ref{sec:floors}.

\subsection{Artificial viscosity}
\label{sec:av}

Shock capturing uses a \emph{frozen} Laplacian artificial viscosity: a fixed
scalar field $a(\bx)\ge 0$ adds isotropic diffusion to all four equations,
\begin{equation}
F_{\mathrm{av}} = a(\bx)\,\nabla\bu ,
\end{equation}
i.e.\ the term $-\nabla\cdot(a\nabla\bu)$ on the left side of \eqref{eq:ns}.
The field $a$ is \emph{data}: it is never differentiated with respect to the
state and never updated inside a Newton solve, so the exact Jacobian is
preserved. Two sources are supported: the archived Exasim field
(\code{vdg.bin}, which equals $0.025\tanh(5(r-1))$ to $2.8\times10^{-17}$) and
the analytic family
\begin{equation}
a(\bx) = \lambda \tanh\!\bigl(c\,d(\bx)\bigr),
\qquad d(\bx) = \sqrt{x^2+y^2} - 1
\end{equation}
($d$ is the wall distance; exact for this geometry). The M4 continuation
lowers $(\lambda, c)$ through the schedule of \S\ref{sec:continuation}.

\subsection{First-order (mixed) form and the gradient sign convention}

The HDG discretization rewrites \eqref{eq:ns} as a first-order system with an
auxiliary gradient variable. Following Exasim's convention \emph{verbatim},
\begin{equation}
\bq \;=\; -\nabla \bu \qquad (\text{eight components in 2D}),
\end{equation}
\begin{equation}
\bq + \nabla\bu = 0, \qquad
\nabla\cdot F(\bu, \bq) = 0 .
\end{equation}
All model formulas are then written in the ``model-signed'' gradient slots.
For example, with $r_x \equiv -\partial_x\rho$ etc., the transcribed
$x$-flux reads
\begin{equation}
f_x = \begin{pmatrix}
\rho u + a\,r_x\\[2pt]
\rho u^2 + p + \tau_{xx} + a\,(ru)_x\\[2pt]
\rho u v + \tau_{xy} + a\,(rv)_x\\[2pt]
\rho u H + u\tau_{xx} + v\tau_{xy} + f_c T_x + a\,(rE)_x
\end{pmatrix},
\label{eq:modelflux}
\end{equation}
where $\tau_{xx},\dots,T_x$ are built from the model-signed slots, so the
\emph{plus} signs in \eqref{eq:modelflux} are dissipative. Adopting this
convention project-wide lets the generated Exasim model code (values and
analytic Jacobians) be transcribed without any sign surgery; a dissipativity
unit test (linearized entropy production $\le 0$) pins the convention.

\subsection{Density and pressure regularization (``floors'')}
\label{sec:floors}

The physics model exists in two variants, selected at run time
(\code{physics.regularization} in the YAML deck):

\paragraph{Floor-free model (default).} Equations
\eqref{eq:ns}--\eqref{eq:modelflux} exactly as exported in the reference
run's \code{pdemodel.txt}. This is the model certified by the M3 parity
gates.

\paragraph{Regularized model (\code{floors}).} The model of the original
\code{pdemodel\_ns.m}, which clamps density and pressure through the smoothed
hinge
\begin{equation}
\mathrm{lmax}(x;\alpha) \;=\;
x\left(\frac{\arctan(\alpha x)}{\pi} + \frac12\right)
- \frac{\arctan\alpha}{\pi} + \frac12,
\qquad \alpha = 10^{3},
\end{equation}
which satisfies $\mathrm{lmax}(x)\to x$ for $x\gg 1/\alpha$,
$\mathrm{lmax}(x)\to 0$ for $x\ll -1/\alpha$. Inside the flux only,
\begin{align}
\rho &\leftarrow \rho_{\min} + \mathrm{lmax}(\rho - \rho_{\min};\alpha),
& \rho_{\min} &= 10^{-2},\\
p &\leftarrow p_{\min} + \mathrm{lmax}(p - p_{\min};\alpha),
& p_{\min} &= 10^{-3},
\end{align}
and the density and pressure gradient slots are scaled by the corresponding
sensors (the derivative of the clamp),
\begin{equation}
d(x) = \frac{\arctan(\alpha \tilde x)}{\pi}
+ \frac{\alpha \tilde x}{\pi(\alpha^2 \tilde x^2 + 1)} + \frac12,
\qquad \tilde x = (\text{clamped value}) - (\text{floor}),
\end{equation}
applied as $(r_x, r_y) \leftarrow d_\rho\,(r_x, r_y)$ and
$(p_x, p_y) \leftarrow d_p\,(p_x, p_y)$. The Sutherland law
\eqref{eq:sutherland} is evaluated at the clamped $(\rho, p)$, so the flux is
a smooth, globally defined function of the state --- in particular finite at
$p<0$.

Two properties make this safe and necessary, respectively:
\begin{itemize}
\item \textbf{Inactive at converged states.} The clamp deviates from the
identity by $O\!\bigl(1/(\alpha^3 x^2)\bigr)$; at the converged states of
this problem ($\min p \approx 0.0111 \gg p_{\min}$) the two models agree to
$\sim\!10^{-9}$ relative. All M3 parity gates are therefore blind to the
difference, and both models share the same fixed points.
\item \textbf{Load-bearing in the continuation.} Lowering the artificial
viscosity from $0.06$ to $0.04$ sharpens the captured shock; the first Newton
correction necessarily undershoots to $p \approx -0.009$ at the
stagnation-line shock foot (both this code and Exasim's own solver produce
this trial state, agreeing to $0.02\%$). Without floors the trial residual is
NaN (via $T_r^{3/2}$) and no line-search or pseudo-transient strategy can
recover, because deciding to damp already requires evaluating the residual.
With floors the same trial evaluates to a small finite residual, the step is
accepted, and Newton restores positivity by itself within three iterations.
The full forensic account is in \code{M4\_FAILURE\_REPORT.md}.
\end{itemize}

The regularized flux and its analytic Jacobian in \code{ns\_physics.hpp} are
mechanical transcriptions of symbolically generated code
(\code{m4\_diagnostic/my\_model\_floored.hpp}, produced by Exasim's
\code{text2code} from the floored model source); a unit test pins the
transcription at exactly zero error over admissible \emph{and}
negative-pressure states.

%======================================================================
\section{Nondimensionalization}
%======================================================================
\label{sec:nondim}

All quantities are scaled by freestream density $\rho_\infty$, freestream
speed $V_\infty$, and the cylinder radius $R$:
\begin{equation}
\rho^* = \frac{\rho}{\rho_\infty},\quad
(u^*,v^*) = \frac{(u,v)}{V_\infty},\quad
p^* = \frac{p}{\rho_\infty V_\infty^2},\quad
E^* = \frac{E}{V_\infty^2},\quad
\bx^* = \frac{\bx}{R},\quad
\mu^* = \frac{\mu}{\rho_\infty V_\infty R}.
\end{equation}
(The asterisks are dropped everywhere else in this document.) Consequently
\begin{equation}
\rho_\infty^* = 1,\qquad
(u,v)_\infty^* = (1, 0),\qquad
p_\infty = \frac{1}{\gamma \Minf^2} = 0.011077\ldots,\qquad
\mu_\infty^* = \frac{1}{\Rey},
\end{equation}
\begin{equation}
T_\infty = \frac{p_\infty}{(\gamma-1)\rho_\infty}
= \frac{1}{\gamma(\gamma-1)\Minf^2} = 0.0276937\ldots,
\qquad
(\rho E)_\infty = \frac12 + \frac{p_\infty}{\gamma-1} = 0.5276937\ldots
\end{equation}
Physical temperature is recovered through
$T_{\mathrm{phys}} = T_{\mathrm{ref}}\, T/T_\infty$ with
$T_{\mathrm{ref}} = 124.49$~K; the nondimensional isothermal-wall value is
\begin{equation}
T_{\mathrm{isoW}} \;=\; T_\infty\,
\frac{T_{\mathrm{wall}}}{T_{\mathrm{ref}}} .
\end{equation}

\subsection{The rounded deck constants (deliberately replicated)}
\label{sec:rounded}

The reference run's parameter vector (in the code:
\code{NSParams::mu[0..10]}) is, verbatim,
\begin{equation*}
[\gamma,\ \Rey,\ \Pran,\ \Minf,\ \rho_\infty,\ (\rho u)_\infty,\
(\rho v)_\infty,\ (\rho E)_\infty,\ T_\infty,\ T_{\mathrm{ref}},\
T_{\mathrm{wall}}]
\end{equation*}
\begin{equation*}
= [1.4,\ 183500,\ 0.71,\ 8.03,\ 1,\ 1,\ 0,\
\mathbf{0.52769},\ \mathbf{0.027694},\ 124.49,\ 294.44].
\end{equation*}
The bold entries are \emph{rounded literals}: the exact values are
$0.5276937\ldots$ and $0.0276937\ldots$, so the deck is internally
inconsistent at the $4\times10^{-6}$ level. The model uses the two versions
in different places, and this split must be replicated exactly (mixing them
up costs $1.1\times10^{-5}$ relative in the fields --- enough to fail the
parity gates):
\begin{itemize}
\item the \textbf{flux/Sutherland path} never reads the rounded
$T_\infty$; it forms $T_r = \gamma\Minf^2 p/\rho$ algebraically, i.e.\ uses
the full-precision $T_\infty = 1/(\gamma(\gamma-1)\Minf^2)$ derived from
$\gamma$ and $\Minf$ (\code{NSParams::TinfFlux()});
\item the \textbf{wall boundary condition} uses the rounded literal
verbatim:
$T_{\mathrm{isoW}} = 0.027694 \times 294.44 / 124.49
= 0.06550101502128686$ (\code{NSParams::TisoW()}).
\end{itemize}
The freestream state used in Newton initializations is likewise the rounded
$(1, 1, 0, 0.52769)$.

%======================================================================
\section{Boundary conditions}
%======================================================================
\label{sec:bcs}

Boundary conditions are imposed \emph{weakly through the trace unknown}
$\buh$ (\S\ref{sec:hdg}): on a boundary face the flux-continuity rows of the
trace equation are replaced by residual rows
$\oint_F \mu\, f_b\,\mathrm{d}s = 0$, with $f_b$ per boundary type
($\hat u_i$ denote trace components, $u_i$ interior-state components):
\begin{center}
\begin{tabular}{llll}
\toprule
Attribute & Boundary & Type & $f_b$ \\
\midrule
1 (21 faces) & cylinder $r=1$ & isothermal no-slip wall &
$\bigl(u_0-\hat u_0,\ -\hat u_1,\ -\hat u_2,\
\hat u_0 T_{\mathrm{isoW}} - \hat u_3\bigr)$\\
2 (62 faces) & $x=0$ plane & supersonic outflow &
$\bu - \buh$ (extrapolation)\\
3 (21 faces) & outer arc & supersonic freestream &
$\bu_\infty - \buh$ (Dirichlet on trace)\\
\bottomrule
\end{tabular}
\end{center}
The wall rows extrapolate density, set zero velocity, and set trace internal
energy to $T_{\mathrm{isoW}}$ (via $\widehat{\rho E} = \hat\rho\,
T_{\mathrm{isoW}}$ at zero velocity). Two structural points, both verified
against the generated Exasim code: the stabilization parameter $\tau$ does
\emph{not} appear in $f_b$; and the \emph{element} equation is unchanged on
boundary faces (it keeps the interior numerical flux $\widehat f$) --- only
the trace rows change. The energy wall row is bilinear in $\buh$, so its
Jacobian carries both $-1$ and $+T_{\mathrm{isoW}}$ entries.

%======================================================================
\section{HDG discretization}
%======================================================================
\label{sec:hdg}

\subsection{Mesh and geometry}

The mesh is the reference run's mesh: 651 curved quadrilaterals
($31\times 21$, graded toward the wall), 704 vertices, 1{,}354 edges (104 on
the boundary). Topology comes from \code{grid.bin}; the curved geometry is
taken \emph{exactly} from \code{xdg.bin}, the discrete $Q_4$ coordinate map
on Exasim's tensor Chebyshev--Gauss--Lobatto nodes (1D nodes
$[0,\ (3-\sqrt5)/4,\ 1/2,\ (1+\sqrt5)/4,\ 1]$). Because a $Q_4$ polynomial
re-expressed in a different nodal basis is exact, the MFEM $H^1$ $Q_4$ mesh
nodes are set by evaluating Exasim's basis at MFEM's node locations; the
resulting geometry reproduces \code{xdg.bin} to $1.1\times10^{-15}$. An
analytic mesh generator (same $31\times21$ grading law:
$\theta = 3\pi/2 - \pi\,\eta_2$, outer radius $4.7 + 1.7\cos\theta$, radial
stretching $\mathrm{logdec}(\eta_1,5) = (1-e^{-5\eta_1})/(1-e^{-5})$)
provides meshes for convergence studies and an independent M4 variant.

\subsection{Spaces and unknowns}

All spaces are \emph{broken} (discontinuous) tensor-product polynomials of
degree $p=4$ on the reference square:
\begin{itemize}
\item $\bu_h$: 4 components $\times$ 25 nodes per element
($651\times100 = 65{,}100$ unknowns);
\item $\bq_h$: 8 components $\times$ 25 nodes per element (never a solver
unknown --- eliminated exactly, see below);
\item $\buh_h$: single-valued polynomial of degree 4 on every mesh edge
(MFEM \code{DG\_Interface\_FECollection}), 4 components $\times$ 5 nodes
$\times$ 1{,}354 faces $= 27{,}080$ unknowns --- \emph{the only globally
coupled unknowns}. Boundary faces carry genuine trace unknowns (BCs are
weak; nothing is eliminated strongly).
\end{itemize}
The trace space is stored with component-fastest ordering per face (20
contiguous scalar dofs per face), matching Exasim's layout and giving clean
$20\times20$ face blocks.

\subsection{Weak forms}

Let $\varphi$ denote element test functions and $\mu$ face test functions.
Per element $K$:

\paragraph{Gradient equation (linear, geometry-only).} For each scalar
component and direction $d$,
\begin{equation}
\int_K \bq\,\varphi
= \int_K \bu\,\partial_{x_d}\varphi
- \oint_{\partial K} \buh\, n_d\, \varphi ,
\qquad\Longrightarrow\qquad
\bq_d = C_d\,\bu - E_d\,\buh
\label{eq:qrel}
\end{equation}
with precomputed $25\times25$ and $25\times20$ matrices
$C_d = M^{-1}\!\int \partial_d\varphi\,\varphi$ and
$E_d = M^{-1}\!\oint \varphi\,\varphi_{\mathrm{face}} n_d$ (shared by all
components; geometry-constant). $\bq$ is eliminated through
\eqref{eq:qrel} from the start.

\paragraph{Element equation.}
\begin{equation}
R_u \;=\; \int_K \nabla\varphi : F(\bu, \bq, a)\,\mathrm{d}\bx
\;-\; \oint_{\partial K} \varphi\, \widehat f\,\mathrm{d}s \;=\; 0,
\end{equation}
with the HDG numerical flux
\begin{equation}
\widehat f \;=\; F(\buh, \bq)\cdot\bn \;+\; \tau\,(\bu - \buh),
\qquad \tau = 1 \ \text{(all equations, all faces)}.
\label{eq:fhat}
\end{equation}
The flux in \eqref{eq:fhat} is the \emph{full} (inviscid + viscous + AV) flux
with the state slot replaced by the trace $\buh$ and the gradient slots kept
from the element interior $\bq$ --- this specific evaluation point is
load-bearing for both stability and Jacobian structure, and holds on boundary
faces too.

\paragraph{Trace equation.} On interior faces, flux continuity
\begin{equation}
\oint_F \mu\,\bigl(\widehat f\big|_{K_1} + \widehat f\big|_{K_2}\bigr)
\mathrm{d}s = 0
\end{equation}
(outward normals opposite); on boundary faces the rows are replaced by the
weak boundary residuals of \S\ref{sec:bcs}, mass-weighted with the same
5-point face rule. Because Exasim assembles the same weighted residuals,
residual \emph{norms} are comparable between the two codes, not just
solutions.

\subsection{Quadrature}

Volume: $5\times5$ tensor Gauss--Legendre (degree 8, 25 points); faces:
5-point Gauss--Legendre. These match Exasim's decoded rules exactly. Shape
tables, geometry factors ($\det J$, $\mathrm{adj}\,J$, face normals times
surface Jacobians) and the frozen AV values at all volume/face quadrature
points are tabulated once in the operator constructor (the mesh never moves).

\subsection{Newton linearization and static condensation}

Each Newton iteration assembles, per element, the exact-Jacobian blocks
\begin{equation}
A = \frac{\partial R_u}{\partial \bu},\quad
B = \frac{\partial R_u}{\partial \bq},\quad
F = \frac{\partial R_u}{\partial \buh},\quad
K = \frac{\partial R_h}{\partial \bu},\quad
G = \frac{\partial R_h}{\partial \bq},\quad
H = \frac{\partial R_h}{\partial \buh},
\end{equation}
folds the gradient dependence through \eqref{eq:qrel}
($A \mathrel{+}= \sum_d B_d C_d$, $F \mathrel{-}= \sum_d B_d E_d$, and
likewise for $K, G, H$), and condenses:
\begin{equation}
H_c = H - K A^{-1} F \ (80\times80),\qquad
b_c = K A^{-1} R_u - R_h .
\end{equation}
Eliminating $\bq$ first reduces the per-element factorization from
$300\times300$ to $100\times100$ (a $27\times$ flop reduction); peak solver
memory is $\approx 106$~MB. The condensed system
$H_c^{\mathrm{glob}}\,\Delta\buh = b_c^{\mathrm{glob}}$ is a
$27{,}080\times27{,}080$ sparse matrix with $\le 140$ nonzeros per row (each
face couples to itself and at most six sibling faces). After the trace solve,
the local recovery is
\begin{equation}
\Delta\bu = -A^{-1}(R_u + F\,\Delta\buh),
\qquad\text{then}\qquad
\bq := C\bu - E\buh \ \ \text{(full recompute)}.
\end{equation}

%======================================================================
\section{Solution algorithm}
%======================================================================

\subsection{Damped Newton (Exasim semantics, replicated exactly)}
\label{sec:newton}

Per iteration: full reassembly at the current $(\bu, \bq, \buh)$; direct
solve of the condensed system; update $\buh \mathrel{+}= \alpha\,\Delta\buh$
with local recovery of $\bu$ and recomputation of $\bq$; then the line
search. The line-search and abort rules are transcribed from Exasim's
\code{NewtonSolver}:
\begin{itemize}
\item trial $\alpha = 1$; \textbf{while} the residual increased
\emph{and} $\alpha > 0.1$: halve $\alpha$ and retry. Because the guard tests
$\alpha$ \emph{before} halving, the admissible steps are
$\alpha \in \{1, \tfrac12, \tfrac14, \tfrac18, \tfrac1{16}\}$ and the step
$\alpha = 0.0625$ is \emph{accepted even if the residual still increased}
--- a quirk that is load-bearing at Mach 8;
\item hard abort only if the trial residual is NaN, or increased
\emph{and} exceeds $10^6$;
\item convergence when $\lVert R_u\rVert + \lVert R_h\rVert < 10^{-6}$
(stacked-norm definition, matching Exasim); at most 30 iterations.
\end{itemize}

\subsection{Linear solver (deliberate deviation)}

The condensed systems are solved by a sparse direct factorization (MUMPS via
PETSc, \code{KSPPREONLY} + \code{PCLU}). Exasim uses restarted GMRES with a
degree-20 polynomial preconditioner; that choice affects inner iteration
counts only, not the Newton fixed point, and is deliberately not replicated.
An optional GMRES + face-block-Jacobi mode exists for solver studies. A
pseudo-transient continuation hook (backward-Euler term $(1/\Delta\tau)$
added to $A$, SER ramp) is wired but \emph{off} in all passing
configurations --- the reference schedule needs no PTC.

\subsection{Initialization modes}

\begin{itemize}
\item \code{udg\_file} (M3): load $\bu$ from Exasim's \code{udg.bin}
(components 0--3 only), set $\buh$ to the one-sided trace of $\bu$ (from the
lower-index element on each face, matching Exasim's restart), then recompute
$\bq = C\bu - E\buh$. The stored gradient slots are deliberately discarded,
exactly as Exasim's solver does on restart.
\item \code{damped\_freestream} (M4): the \code{pdeapp\_ns.m} initial
condition
\begin{equation}
\rho = 1,\quad
\rho u = \tanh(10\,d),\quad
\rho v = 0,\quad
\rho E = T_\infty\!\left[\Bigl(\tfrac{T_{\mathrm{wall}}}{T_{\mathrm{ref}}}
- 1\Bigr)e^{-10 d} + 1\right] + \tfrac12\tanh^2(10\,d),
\end{equation}
with $d = r - 1$ the wall distance: freestream flow smoothly brought to rest
at the wall with a near-wall thermal layer consistent with the hot wall.
\end{itemize}

\subsection{Artificial-viscosity continuation (M4)}
\label{sec:continuation}

The self-contained solve replicates the reference script's four steady
solves with frozen AV fields, each warm-starting from the previous:
\begin{equation}
a_k(\bx) = \lambda_k \tanh(c_k\, d),\qquad
(\lambda_k, c_k) \in \bigl\{(0.06, 30),\ (0.04, 30),\
(0.025, 30),\ (0.025, 5)\bigr\}.
\end{equation}
At each stage start the trace is re-initialized to the one-sided interior
trace and $\bq$ recomputed (Exasim's per-stage restart semantics; this makes
the per-stage residual histories directly comparable between the codes). The
regularized model of \S\ref{sec:floors} is required from stage 2 on; with it,
the run needs no PTC, no extra stages, and no damping after stage 1:
\begin{center}
\begin{tabular}{clccc}
\toprule
Stage & AV field & Newton updates & Damped steps & Final residual\\
\midrule
1 & $0.06\tanh(30 d)$ & 9 & 5 (min $\alpha = 0.125$) & $2.52\times10^{-9}$\\
2 & $0.04\tanh(30 d)$ & 4 & 0 & $2.30\times10^{-8}$\\
3 & $0.025\tanh(30 d)$ & 4 & 0 & $2.49\times10^{-7}$\\
4 & $0.025\tanh(5 d)$ & 4 & 0 & $3.30\times10^{-7}$\\
\bottomrule
\end{tabular}
\end{center}
The $9/4/4/4$ signature matches Exasim's own (floored) runs of the same
schedule; the stage-2 histories track step for step, both accepting a first
trial whose minimum pressure is $\approx -0.0090$ (absorbed by the floors).
The end state matches the M3 replication to $3.4\times10^{-6}$ in wall heat
flux --- far inside the 1\% acceptance gate.

\subsection{Wall outputs}

Wall postprocessing evaluates, at the 5 Gauss points of each wall face and
with the trace state in the state slot and interior $\bq$ in the gradient
slots (exactly the reference \code{Fint} functional),
\begin{equation}
q_w = f_c\,(T_x n_x + T_y n_y) + \tau\,(\rho E - \widehat{\rho E}),
\qquad
C_p = 2\,(p(\buh) - p_\infty),
\end{equation}
the latter because the dynamic pressure is
$\tfrac12\rho_\infty V_\infty^2 = \tfrac12$ in these units. The shock
standoff is located as the maximum of $|\partial_x \rho|$ along the
stagnation line from 1{,}000 samples.

%======================================================================
\section{Code implementation}
%======================================================================
\label{sec:code}

\subsection{Layout}

The app is a flat directory of small translation units over MFEM 4.8.1-dev
(MPI + METIS + PETSc/MUMPS build; serial execution --- 651 elements is
decisively serial territory):
\begin{center}
\begin{tabular}{lp{0.66\textwidth}}
\toprule
File & Contents\\
\midrule
\code{ns\_physics.hpp} & Header-only physics on raw \code{double} arrays (no
MFEM types; shared with the planned DG driver): fluxes, Sutherland, boundary
residuals, analytic Jacobians transcribed from the generated Exasim model;
the regularized (floored) flux/Jacobian pair behind
\code{NSParams::regularized}; sign-flip adapters for physical-gradient
callers.\\
\code{exasim\_io} & Exasim binary readers (3-double header + column-major
payload), Chebyshev--Gauss--Lobatto basis evaluation, exact $25\times25$
change-of-basis, per-rank \code{elempart} parsing for the partitioned
reference outputs.\\
\code{exasim\_mesh} & Mesh converter (\code{grid.bin} topology +
\code{xdg.bin} exact geometry + boundary attribution by the run's own
predicates) and the analytic mesh generator.\\
\code{hdg\_ns\_operator} & The HDG core: precomputed tables, $C_d/E_d$,
residual and Jacobian assembly, boundary-row substitution, static
condensation into an MFEM \code{SparseMatrix}, local recovery,
residual-only path, diagnostics (min $\rho$/$p$, symmetry error).\\
\code{hdg\_newton.hpp} & Damped Newton of \S\ref{sec:newton} with the PETSc
direct solve and the optional PTC hook; failures are reported, not thrown.\\
\code{wall\_post} & $C_p$, \code{Fint} heat flux, shock standoff, and the
field/wall comparison machinery used by the M3/M4 gates.\\
\code{hdg\_ns\_driver.cpp} & YAML-configured driver (mesh source, physics,
AV mode, init mode, Newton/PTC, continuation schedule, outputs); ParaView
output of conservative/primitive/AV fields.\\
\code{test\_*} & The verification suite (\S\ref{sec:verification}).\\
\code{m4\_diagnostic/} & Forensic artifacts of the M4 investigation: the
floored model source and its generated code (the transcription reference),
stage-driving script, logs, and the first-failing-step dumps.\\
\bottomrule
\end{tabular}
\end{center}

\subsection{Selected implementation details}

\begin{itemize}
\item \textbf{Exact geometry and state transfer.} Both the $Q_4$ geometry
and all state I/O go through exact polynomial change-of-basis between
Exasim's Chebyshev--Lobatto nodes and MFEM's node sets (round-trip error
$\le 7\times10^{-14}$); naive nodal copying would be $O(10^{-3})$ wrong at
$p=4$.
\item \textbf{Jacobians are transcribed, not derived.} Both model variants'
analytic Jacobians come from Exasim's symbolic code generator; unit tests
compare the transcriptions against lightly adapted copies of the generated
headers at $10^{-13}$ (the floored pair matches \emph{exactly}), plus
central finite-difference checks. Hand-differentiation of the Sutherland/
sensor chain rules was deliberately avoided.
\item \textbf{Sparsity with explicit zeros.} The condensed matrix graph is
built once with \code{skip\_zeros=0}; exact zeros are guaranteed in boundary
blocks and dropping them would corrupt later re-assemblies.
\item \textbf{Frozen-AV tabulation.} $a(\bx)$ is tabulated once at all
volume and face quadrature points; it enters the Jacobian only as
$\partial f_{c,d}/\partial q_{c,d} \mathrel{+}= a$ (identity blocks scaled
by data).
\item \textbf{Face transformations.} Face physical coordinates are always
evaluated through the face transformation objects; MFEM reuses mutable
transformation storage, and retaining stale element pointers across face
queries corrupts face data (found the hard way in M2).
\item \textbf{Build hygiene.} Every translation unit that includes
\code{ns\_physics.hpp} lists it as a makefile dependency; \code{NSParams}
is passed across TU boundaries, and a stale-object ABI mismatch after the
struct grew produced a heap abort before the dependencies were fixed.
\item \textbf{Timings} (single process, development host): a full M3 Newton
solve (4 iterations) spends $\approx 5.1$~s in assembly and $\approx 1.1$~s
in the direct trace solves; each M4 stage runs in $\approx 7$--$14$~s.
\end{itemize}

\subsection{Verification gates}
\label{sec:verification}

\code{make test} runs the milestone gates in sequence:
\begin{center}
\small
\begin{tabular}{lp{0.5\textwidth}l}
\toprule
Gate & Statement & Result\\
\midrule
M1 physics & flux/BC values and all Jacobians vs the generated model at
1000 random states; FD consistency; dissipativity sign & $\le 10^{-13}$;
$4\times10^{-9}$\\
M1 regularized & floored flux/Jacobian vs generated floored model over
admissible + floor-straddling (incl.\ $p<0$) states; FD; floored/unfloored
agreement at admissible states; M4 failure state finite (floored) and NaN
(unfloored) & exactly 0; $4.9\times10^{-7}$; $4.2\times10^{-7}$; pass\\
M1 mesh/IO & topology/attributes; exact \code{xdg} geometry; change-of-basis
round trips; quadrature identity & $1.1\times10^{-15}$\\
M2 core & trace-orientation consistency; global condensed Jacobian vs
directional FD; MMS convergence order $\ge 4.8$; freestream preservation;
$M{=}4.5$ sanity case & pass ($9.7\times10^{-10}$ FD)\\
M3 replication & residual of Exasim's converged state in this assembler;
Newton from \code{udg.bin}; field $L^2$ vs Exasim; wall $C_p$/heat
flux/standoff & $3.3\times10^{-7}$; $\sim10^{-10}$;
$\sim10^{-9}$\\
M4 continuation & four-stage schedule (no PTC) on converted and analytic
meshes; final wall heat flux and standoff vs M3 within 1\% & $9/4/4/4$;
$3.4\times10^{-6}$\\
\bottomrule
\end{tabular}
\end{center}
The strongest single test is the M3 residual cross-check: Exasim's converged
$(\bu, \bq, \buh)$, loaded into this code's assembler, must evaluate to a
residual at Exasim's own convergence tolerance. It certifies discretization
parity independent of any solver behavior, and it gates everything
downstream.

\subsection{Running}

\begin{verbatim}
make            # builds driver + all test binaries
make test       # M1-M4 verification gates
./hdg_ns_driver -i Input/input_mach8_baseline.yaml       # M3 replication
./hdg_ns_driver -i Input/input_mach8_continuation.yaml   # M4 continuation
\end{verbatim}
A deck selects the mesh source (\code{exasim} converter or \code{analytic}),
physics constants and regularization, AV mode (\code{file} or \code{tanh}),
initialization, Newton/PTC settings, the continuation schedule, and output
paths. Outputs are high-order ParaView collections
(conservative/primitive/AV fields), wall CSVs
$(\theta, x, y, C_p, q_w)$ at wall Gauss points, and per-stage
continuation history/summary CSVs.

%======================================================================
\section*{Further reading in the repository}
%======================================================================

\begin{itemize}
\item \code{PLAN\_HDG.md} --- the full implementation plan, with
file/line references into the Exasim sources for every discretization
decision.
\item \code{M4\_FAILURE\_REPORT.md} --- the continuation failure forensics:
the negative-pressure mechanism, the cross-code first-failing-step
comparison, and the floored-model reconstruction that resolved it.
\item \code{NOTES.md} --- milestone-to-milestone findings (data-format
gotchas, orderings, measured parity numbers).
\item \code{PLAN\_DG.md} --- the planned sibling nodal-DG driver sharing
\code{ns\_physics}/\code{exasim\_io}/\code{exasim\_mesh}.
\end{itemize}

\end{document}
